\phantomsection
\addcontentsline{toc}{section}{References}
\begin{thebibliography}{99}
\bibitem{bcw} J.~Bausch, T.~S.~Cubitt and J.~D.~Watson.
Uncomputability of phase diagrams.
\emph{Nature Communications} 12 (2021), 452.
\href{https://doi.org/10.1038/s41467-020-20504-6}{doi:10.1038/s41467-020-20504-6}.

\bibitem{bmm} A.~Bu\v{s}i\'c, J.~Mairesse and I.~Marcovici.
Probabilistic cellular automata, invariant measures, and perfect sampling.
\emph{Advances in Applied Probability} 45 (2013), 960--980.
\href{https://doi.org/10.1239/aap/1386857853}{doi:10.1239/aap/1386857853}.

\bibitem{cpw} T.~S.~Cubitt, D.~P\'erez-Garc\'ia and M.~M.~Wolf.
Undecidability of the spectral gap.
\emph{Nature} 528 (2015), 207--211.
\href{https://doi.org/10.1038/nature16059}{doi:10.1038/nature16059}.

\bibitem{dk} E.~Domany and W.~Kinzel.
Equivalence of cellular automata to Ising models and directed percolation.
\emph{Physical Review Letters} 53 (1984), 311--314.
\href{https://doi.org/10.1103/PhysRevLett.53.311}{doi:10.1103/PhysRevLett.53.311}.

\bibitem{dmForgetting} S.~Douglas and N.~Mghirbi.
Forgetting under noise is undecidable: ergodicity of positive-rate cellular automata at a fixed noise rate.
Preprint, version 3.0.0 (2026).
\href{https://doi.org/10.5281/zenodo.24265098}{doi:10.5281/zenodo.24265098}.

\bibitem{dmNoise} S.~Douglas and N.~Mghirbi.
One noisy cellular automaton with an uncomputable noise diagram.
Preprint, version 1.4.0 (2026).
\href{https://doi.org/10.5281/zenodo.24204195}{doi:10.5281/zenodo.24204195}.

\bibitem{dmWindows} S.~Douglas and N.~Mghirbi.
A perturbed cellular automaton with infinitely many ergodicity transitions.
Preprint, version 1.0.0 (2026).
\href{https://doi.org/10.5281/zenodo.24122609}{doi:10.5281/zenodo.24122609}.

\bibitem{dmSpaceTime} S.~Douglas and N.~Mghirbi.
Exponential ergodicity and the space-time Gibbs states of one-dimensional probabilistic cellular automata.
Preprint, version 1.2.0 (2026).
\href{https://doi.org/10.5281/zenodo.24122307}{doi:10.5281/zenodo.24122307}.

\bibitem{dmHard} S.~Douglas and N.~Mghirbi.
An uncomputable phase diagram at positive temperature for one classical lattice model.
Preprint, version 1.2.0 (2026).
\href{https://doi.org/10.5281/zenodo.24144192}{doi:10.5281/zenodo.24144192}.

\bibitem{gacs86} P.~G\'acs.
Reliable computation with cellular automata.
\emph{Journal of Computer and System Sciences} 32 (1986), 15--78.
\href{https://doi.org/10.1016/0022-0000(86)90002-4}{doi:10.1016/0022-0000(86)90002-4}.

\bibitem{gacs01} P.~G\'acs.
Reliable cellular automata with self-organization.
\emph{Journal of Statistical Physics} 103 (2001), 45--267.
Corrected and strengthened version \href{https://arxiv.org/abs/math/0003117v2}{arXiv:math/0003117v2}.

\bibitem{gklm} S.~Goldstein, R.~Kuik, J.~L.~Lebowitz and C.~Maes.
From PCA's to equilibrium systems and back.
\emph{Communications in Mathematical Physics} 125 (1989), 71--79.

\bibitem{kari} J.~Kari.
The nilpotency problem of one-dimensional cellular automata.
\emph{SIAM Journal on Computing} 21 (1992), 571--586.
\href{https://doi.org/10.1137/0221036}{doi:10.1137/0221036}.

\bibitem{lms} J.~L.~Lebowitz, C.~Maes and E.~R.~Speer.
Statistical mechanics of probabilistic cellular automata.
\emph{Journal of Statistical Physics} 59 (1990), 117--170.
\href{https://doi.org/10.1007/BF01015566}{doi:10.1007/BF01015566}.

\bibitem{mtBernoulli} I.~Marcovici and S.~Taati.
Positive-rate PCA and IPS with stationary Bernoulli measures are rapidly forgetful.
Preprint (2026), \href{https://arxiv.org/abs/2605.16904}{arXiv:2605.16904}.

\bibitem{msStability} H.~Marsan and M.~Sablik.
Strong stochastic stability of cellular automata.
Preprint (2024), \href{https://arxiv.org/abs/2403.19932}{arXiv:2403.19932}.

\bibitem{msMaps} H.~Marsan and M.~Sablik.
Flexibility versus genericity of phase diagrams of perturbed continuous maps on the Cantor set.
Preprint (2025), \href{https://arxiv.org/abs/2508.00461}{arXiv:2508.00461}.

\bibitem{mst} H.~Marsan, M.~Sablik and I.~T\"orm\"a.
A perturbed cellular automaton with two phase transitions for the ergodicity.
\emph{Journal of Statistical Physics} 193 (2026), article 6.
\href{https://doi.org/10.1007/s10955-025-03564-0}{doi:10.1007/s10955-025-03564-0}.

\bibitem{mtaati} J.~Mi\k{e}kisz and S.~Taati.
Quasicrystalline Gibbs states in 4-dimensional lattice-gas models with finite-range interactions.
Preprint (2026), \href{https://arxiv.org/abs/2512.24436v2}{arXiv:2512.24436v2}.

\bibitem{toom} A.~L.~Toom.
Algorithmical unsolvability of the ergodicity problem for binary cellular automata.
\emph{Markov Processes and Related Fields} 6 (2000), 569--577.

\bibitem{toomChapter} A.~L.~Toom, N.~B.~Vasilyev, O.~N.~Stavskaya, L.~G.~Mityushin,
G.~L.~Kurdyumov and S.~A.~Pirogov.
Discrete local Markov systems.
In R.~L.~Dobrushin, V.~I.~Kryukov and A.~L.~Toom (eds.),
\emph{Stochastic Cellular Systems: Ergodicity, Memory, Morphogenesis},
Manchester University Press, 1990, pp.~1--182.

\bibitem{gwpn} M.~Gu, C.~Weedbrook, A.~Perales and M.~A.~Nielsen.
More really is different.
\emph{Physica D} 238 (2009), 835--839.
\href{https://doi.org/10.1016/j.physd.2008.12.016}{doi:10.1016/j.physd.2008.12.016}.

\bibitem{ch} J.-R.~Chazottes and M.~Hochman.
On the zero-temperature limit of Gibbs states.
\emph{Communications in Mathematical Physics} 297 (2010), 265--281.
\href{https://doi.org/10.1007/s00220-010-0997-8}{doi:10.1007/s00220-010-0997-8}.

\bibitem{cs} J.-R.~Chazottes and M.~Shinoda.
On the absence of zero-temperature limit of equilibrium states for finite-range interactions on the lattice \(\Z^2\).
Preprint (2020), \href{https://arxiv.org/abs/2010.08998v2}{arXiv:2010.08998v2}.

\bibitem{bbdt} S.~Barbieri, R.~Bissacot, G.~Dalle Vedove and P.~Thieullen.
Zero-temperature chaos in bidimensional models with finite-range potentials.
\emph{Advances in Mathematics} 457 (2024), 109906.
\href{https://doi.org/10.1016/j.aim.2024.109906}{doi:10.1016/j.aim.2024.109906}.

\bibitem{gst} L.~Gayral, M.~Sablik and S.~Taati.
Characterisation of the set of ground states of uniformly chaotic finite-range lattice models.
\emph{Nonlinearity} 38 (2025), 075011.
\href{https://doi.org/10.1088/1361-6544/addbb8}{doi:10.1088/1361-6544/addbb8}.

\bibitem{dmMagnetization} S.~Douglas and N.~Mghirbi.
Spontaneous magnetization and thermal decay are undecidable at positive temperature.
Preprint, version 1.2.0 (2026).
\href{https://doi.org/10.5281/zenodo.23962706}{doi:10.5281/zenodo.23962706}.

\bibitem{georgii} H.-O.~Georgii.
\emph{Gibbs Measures and Phase Transitions}, second extended edition.
De Gruyter, 2011.
\href{https://doi.org/10.1515/9783110250329}{doi:10.1515/9783110250329}.


\bibitem{mos} F.~Martinelli, E.~Olivieri and R.~H.~Schonmann.
For 2-D lattice spin systems weak mixing implies strong mixing.
\emph{Communications in Mathematical Physics} 165 (1994), 33--47.
\href{https://doi.org/10.1007/BF02099735}{doi:10.1007/BF02099735}.

\bibitem{spinka} Y.~Spinka.
Finitary codings for spatial mixing Markov random fields.
\emph{Annals of Probability} 48 (2020), 1557--1591.
\href{https://doi.org/10.1214/19-AOP1405}{doi:10.1214/19-AOP1405};
\href{https://arxiv.org/abs/1803.10578}{arXiv:1803.10578}.

\bibitem{syozi} I.~Sy\^ozi.
Statistics of Kagom\'e lattice.
\emph{Progress of Theoretical Physics} 6 (1951), 306--308.
\href{https://doi.org/10.1143/ptp/6.3.306}{doi:10.1143/ptp/6.3.306}.

\bibitem{fisher} M.~E.~Fisher.
Transformations of Ising models.
\emph{Physical Review} 113 (1959), 969--981.
\href{https://doi.org/10.1103/PhysRev.113.969}{doi:10.1103/PhysRev.113.969}.

\bibitem{canovaStrecka} L.~\v{C}anov\'a and J.~Stre\v{c}ka.
Reentrant phenomenon in the exactly solvable mixed spin-1/2 and spin-1
Ising--Heisenberg model on diamond-like decorated planar lattices.
\emph{Physica Status Solidi B} 247 (2010), 433--443.
\href{https://doi.org/10.1002/pssb.200945444}{doi:10.1002/pssb.200945444};
\href{https://arxiv.org/abs/0909.4191}{arXiv:0909.4191}.

\bibitem{cencarikovaStreckaLyra} H.~\v{C}en\v{c}arikov\'a, J.~Stre\v{c}ka and M.~L.~Lyra.
Reentrant phase transitions of a coupled spin-electron model on doubly decorated
planar lattices with two or three consecutive critical points.
\emph{Journal of Magnetism and Magnetic Materials} 401 (2016), 1106--1122.
\href{https://doi.org/10.1016/j.jmmm.2015.11.018}{doi:10.1016/j.jmmm.2015.11.018};
\href{https://arxiv.org/abs/1706.09871}{arXiv:1706.09871}.

\bibitem{powersReznick} V.~Powers and B.~Reznick.
A new bound for P\'olya's theorem with applications to polynomials positive on polyhedra.
\emph{Journal of Pure and Applied Algebra} 164 (2001), 221--229.
\href{https://doi.org/10.1016/S0022-4049(00)00155-9}{doi:10.1016/S0022-4049(00)00155-9}.

\bibitem{drt} H.~Duminil-Copin, A.~Raoufi and V.~Tassion.
Sharp phase transition for the random-cluster and Potts models via decision trees.
\emph{Annals of Mathematics} 189 (2019), 75--99.
\href{https://doi.org/10.4007/annals.2019.189.1.2}{doi:10.4007/annals.2019.189.1.2};
\href{https://arxiv.org/abs/1705.03104v2}{arXiv:1705.03104v2}.

\bibitem{dghmt} H.~Duminil-Copin, M.~Gagnebin, M.~Harel, I.~Manolescu and V.~Tassion.
Discontinuity of the phase transition for the planar random-cluster and Potts models with \(q>4\).
\emph{Annales Scientifiques de l'\'Ecole Normale Sup\'erieure} 54 (2021), 1363--1413.
\href{https://doi.org/10.24033/asens.2485}{doi:10.24033/asens.2485};
\href{https://arxiv.org/abs/1611.09877v2}{arXiv:1611.09877v2}.

\bibitem{shapiro} V.~Shapiro.
Semi-analytic geometry with R-functions.
\emph{Acta Numerica} 16 (2007), 239--303.
\href{https://doi.org/10.1017/S096249290631001X}{doi:10.1017/S096249290631001X}.

\bibitem{protasovVoynov} V.~Yu.~Protasov and A.~S.~Voynov.
Matrix semigroups with constant spectral radius.
\emph{Linear Algebra and its Applications} 513 (2017), 376--408.
\href{https://doi.org/10.1016/j.laa.2016.10.013}{doi:10.1016/j.laa.2016.10.013};
\href{https://arxiv.org/abs/1407.6568}{arXiv:1407.6568}.

\bibitem{dsConstructive} R.~L.~Dobrushin and S.~B.~Shlosman.
Constructive criterion for the uniqueness of Gibbs field.
In \emph{Statistical Physics and Dynamical Systems}, Birkh\"auser, Boston (1985), 347--370.
\href{https://doi.org/10.1007/978-1-4899-6653-7_20}{doi:10.1007/978-1-4899-6653-7\_20}.

\bibitem{weitz} D.~Weitz.
Combinatorial criteria for uniqueness of Gibbs measures.
\emph{Random Structures and Algorithms} 27 (2005), 445--475.
\href{https://doi.org/10.1002/rsa.20073}{doi:10.1002/rsa.20073}.

\bibitem{shlosmanDS} S.~Shlosman.
The critical temperature \(T_{cr}\)(Ising) is DS-computable.
Preprint (2025). \href{https://arxiv.org/abs/2505.24750}{arXiv:2505.24750}.

\bibitem{plc} J.~Purcell, Z.~Li and T.~S.~Cubitt.
Chaitin phase transition.
Preprint (2024). \href{https://arxiv.org/abs/2410.02600}{arXiv:2410.02600}.

\bibitem{zarubinKassanOgly} A.~V.~Zarubin and F.~A.~Kassan-Ogly.
The formation of magnetic reentrancy in the Ising model on a decorated square lattice.
Preprint (2026). \href{https://arxiv.org/abs/2606.12572}{arXiv:2606.12572}.

\bibitem{bdc} V.~Beffara and H.~Duminil-Copin.
The self-dual point of the two-dimensional random-cluster model is critical for \(q\ge1\).
\emph{Probability Theory and Related Fields} 153 (2012), 511--542.
\href{https://doi.org/10.1007/s00440-011-0353-8}{doi:10.1007/s00440-011-0353-8}.

\bibitem{stephenson} J.~Stephenson.
Ising model with antiferromagnetic next-nearest-neighbor coupling: spin correlations and disorder points.
\emph{Physical Review B} 1 (1970), 4405--4409.
\href{https://doi.org/10.1103/PhysRevB.1.4405}{doi:10.1103/PhysRevB.1.4405}.

\bibitem{fisherWidom} M.~E.~Fisher and B.~Widom.
Decay of correlations in linear systems.
\emph{Journal of Chemical Physics} 50 (1969), 3756--3772.
\href{https://doi.org/10.1063/1.1671624}{doi:10.1063/1.1671624}.

\bibitem{enting} I.~G.~Enting.
Crystal growth models and Ising models: disorder points.
\emph{Journal of Physics C: Solid State Physics} 10 (1977), 1379--1388.
\href{https://doi.org/10.1088/0022-3719/10/9/011}{doi:10.1088/0022-3719/10/9/011}.

\bibitem{peschelEmery} I.~Peschel and V.~J.~Emery.
Calculation of spin correlations in two-dimensional Ising systems from one-dimensional kinetic models.
\emph{Zeitschrift f\"ur Physik B} 43 (1981), 241--249.
\href{https://doi.org/10.1007/BF01297524}{doi:10.1007/BF01297524}.

\bibitem{rujan} P.~Ruj\'an.
Order and disorder lines in systems with competing interactions: II. The IRF model.
\emph{Journal of Statistical Physics} 29 (1982), 247--262.
\href{https://doi.org/10.1007/BF01020785}{doi:10.1007/BF01020785}.

\bibitem{mairesseMarcovici} J.~Mairesse and I.~Marcovici.
Probabilistic cellular automata and random fields with i.i.d. directions.
\emph{Annales de l'Institut Henri Poincar\'e, Probabilit\'es et Statistiques} 50 (2014), 455--475.
\href{https://doi.org/10.1214/12-AIHP530}{doi:10.1214/12-AIHP530}.

\bibitem{ouaknineWorrell} J.~Ouaknine and J.~Worrell.
Decision problems for linear recurrence sequences.
In \emph{Reachability Problems}, Lecture Notes in Computer Science 7550 (2012), 21--28.
\href{https://doi.org/10.1007/978-3-642-33512-9_3}{doi:10.1007/978-3-642-33512-9\_3}.

\bibitem{caputoParisi} P.~Caputo and D.~Parisi.
Block factorization of the relative entropy via spatial mixing.
\emph{Communications in Mathematical Physics} 388 (2021), 793--818.
\href{https://doi.org/10.1007/s00220-021-04237-1}{doi:10.1007/s00220-021-04237-1}.

\bibitem{casseMarckert} J.~Casse and J.-F.~Marckert.
Markovianity of the invariant distribution of probabilistic cellular automata on the line.
\emph{Stochastic Processes and their Applications} 125 (2015), 3458--3483.
\href{https://doi.org/10.1016/j.spa.2015.05.001}{doi:10.1016/j.spa.2015.05.001};
\href{https://arxiv.org/abs/1401.5365v2}{arXiv:1401.5365v2}.

\bibitem{casseMarcovici} J.~Casse and I.~Marcovici.
Probabilistic cellular automata with memory two: invariant laws and multidirectional reversibility.
\emph{Annales Henri Lebesgue} 3 (2020), 501--559.
\href{https://doi.org/10.5802/ahl.39}{doi:10.5802/ahl.39};
\href{https://arxiv.org/abs/1710.05490}{arXiv:1710.05490}.

\bibitem{dudenhefner} A.~Dudenhefner.
Certified decision procedures for two-counter machines.
In \emph{FSCD 2022}, LIPIcs 228 (2022), 16:1--16:18.
\href{https://doi.org/10.4230/LIPIcs.FSCD.2022.16}{doi:10.4230/LIPIcs.FSCD.2022.16}.

\bibitem{rebeschiniVanHandel} P.~Rebeschini and R.~van Handel.
Comparison theorems for Gibbs measures.
\emph{Journal of Statistical Physics} 157 (2014), 234--281.
\href{https://arxiv.org/abs/1308.4117}{arXiv:1308.4117}.

\bibitem{flajoletSedgewick} P.~Flajolet and R.~Sedgewick.
\emph{Analytic Combinatorics}.
Cambridge University Press, 2009.
\href{https://doi.org/10.1017/CBO9780511801655}{doi:10.1017/CBO9780511801655}.

\end{thebibliography}
